\chapter{Operating Systems, Concurrency and Networks}\label{ch:iv3:operating-systems-concurrency-and-networks}

\begin{ivtrains}
What a process, a thread and a page really are, how a lock reaches the kernel, and what happens to a packet on its way to your code. Everything here refers to Linux on x86-64. Network programs run on loopback, in Python or C++20. Know which behaviour belongs to the protocol, which to Linux, and which to the one run in front of you.
\end{ivtrains}

\begin{question}[$\star$ \iqroles{QD, Dev, LLE} \iqfirm{any} \iqlevel{junior}]\label{iq:iv3:operating-systems-concurrency-and-networks:1}
On Linux, what two lines does \cref{lst:iv3:os:forkthread} print? What does a thread share with the others in its process, and what does each have for itself? What is shared after a \texttt{fork}?
\omcode{code/interviews-3/05-operating-systems-concurrency-and-networks/python/fork_vs_thread.py}{1}{22}{A child process and a thread.\protect\hfil\protect\break}\label{lst:iv3:os:forkthread}
\end{question}

\begin{question}[$\star$ \iqroles{QD, Dev, LLE} \iqfirm{any} \iqlevel{junior}]\label{iq:iv3:operating-systems-concurrency-and-networks:2}
What does \cref{lst:iv3:os:tls} print? Where does each thread's copy of \texttt{mine} live, compared with \texttt{shared}? Why does the last number come out as it does? Name one use of thread-local storage in a counter that is read rarely.
\omcode{code/interviews-3/05-operating-systems-concurrency-and-networks/cpp/tls.cpp}{1}{22}{One variable per thread.\protect\hfil\protect\break}\label{lst:iv3:os:tls}
\end{question}

\begin{question}[$\star$ \iqroles{QD, Dev, LLE} \iqfirm{any} \iqlevel{junior}]\label{iq:iv3:operating-systems-concurrency-and-networks:3}
What three lines does \cref{lst:iv3:os:zombie} print? What is a zombie, what does it still occupy, and how do you get rid of one? What happens to it if the parent dies first?
\omcode{code/interviews-3/05-operating-systems-concurrency-and-networks/python/zombie.py}{1}{18}{A child that has exited.\protect\hfil\protect\break}\label{lst:iv3:os:zombie}
\end{question}

\begin{question}[$\star$ \iqroles{QD, Dev, LLE} \iqfirm{any} \iqlevel{junior}]\label{iq:iv3:operating-systems-concurrency-and-networks:4}
What two lines does \cref{lst:iv3:os:boundaries} print? How many \texttt{recv} calls does the TCP loop need before it ends, and is that number fixed? State what TCP guarantees and UDP does not, and the one thing UDP keeps that TCP does not.
\omcode{code/interviews-3/05-operating-systems-concurrency-and-networks/python/boundaries.py}{1}{19}{Two sends over each protocol.\protect\hfil\protect\break}\label{lst:iv3:os:boundaries}
\end{question}

\begin{question}[$\star$ \iqroles{QD, Dev, LLE} \iqfirm{any} \iqlevel{junior}]\label{iq:iv3:operating-systems-concurrency-and-networks:5}
On x86-64, what three lines does \cref{lst:iv3:os:endian} print? What is network byte order? A feed sends a 32-bit price in little-endian order, so which of the two formats reads it directly on this machine, and what do you need on a big-endian one?
\omcode{code/interviews-3/05-operating-systems-concurrency-and-networks/python/endian.py}{1}{8}{Byte order.\protect\hfil\protect\break}\label{lst:iv3:os:endian}
\end{question}

\begin{question}[$\star$ \iqroles{QD, Dev, LLE} \iqfirm{any} \iqlevel{junior}]\label{iq:iv3:operating-systems-concurrency-and-networks:6}
What three lines does \cref{lst:iv3:os:fivetuple} print? Two clients are connected to one server port, so how does the kernel know which connection a segment belongs to? What identifies a TCP connection?
\omcode{code/interviews-3/05-operating-systems-concurrency-and-networks/python/five_tuple.py}{1}{12}{Two connections to one port.\protect\hfil\protect\break}\label{lst:iv3:os:fivetuple}
\end{question}

\begin{question}[$\star$ \iqroles{QD, Dev, LLE} \iqfirm{any} \iqlevel{junior}]\label{iq:iv3:operating-systems-concurrency-and-networks:7}
With g++ 11.4, what two lines does \cref{lst:iv3:os:semaphore} print? What is the difference between a semaphore and a mutex, given that \texttt{other} released a slot it never acquired? What would that same act be on a \texttt{std::mutex}?
\omcode{code/interviews-3/05-operating-systems-concurrency-and-networks/cpp/semaphore.cpp}{1}{13}{Three slots.\protect\hfil\protect\break}\label{lst:iv3:os:semaphore}
\end{question}

\begin{question}[$\star$ \iqroles{QD, Dev, LLE} \iqfirm{any} \iqlevel{junior}]\label{iq:iv3:operating-systems-concurrency-and-networks:8}
What three lines does \cref{lst:iv3:os:signals} print? Why may a process catch \texttt{SIGTERM} but not \texttt{SIGKILL}? So how should you stop a trading process that holds open orders?
\omcode{code/interviews-3/05-operating-systems-concurrency-and-networks/python/signals.py}{1}{18}{Two signals.\protect\hfil\protect\break}\label{lst:iv3:os:signals}
\end{question}

\begin{question}[$\star\star$ \iqroles{QD, Dev, LLE} \iqfirm{any} \iqlevel{junior, experienced}]\label{iq:iv3:operating-systems-concurrency-and-networks:9}
\Cref{lst:iv3:os:cow} touches 2000 pages of 4 KiB, then forks. Roughly how many page faults does the child take on reading them all, and on then writing them all? What values does each process then print, and how can both print the same address? Explain copy-on-write, and what it means for a large process that forks to take a snapshot while it keeps updating its data.
\omcode{code/interviews-3/05-operating-systems-concurrency-and-networks/cpp/cow.cpp}{8}{37}{Reading and writing after a fork.\protect\hfil\protect\break}\label{lst:iv3:os:cow}
\end{question}

\begin{question}[$\star\star$ \iqroles{QD, Dev, LLE} \iqfirm{any} \iqlevel{junior, experienced}]\label{iq:iv3:operating-systems-concurrency-and-networks:10}
x86-64 with 4-level paging and 4 KiB pages translates 48-bit virtual addresses, with 512 entries of 8 bytes in each table. How is the address \texttt{0x7f1234567abc} split, and what are the four table indices and the offset, in decimal? How much memory do the page tables take to map 1 GiB of contiguous, aligned virtual memory? What changes with 2 MiB pages? Write it in Python.
\end{question}

\begin{question}[$\star\star$ \iqroles{QD, Dev, LLE} \iqfirm{any} \iqlevel{experienced}]\label{iq:iv3:operating-systems-concurrency-and-networks:11}
With glibc on x86-64 Linux, what three lines does \cref{lst:iv3:os:mallocmmap} print, given that \texttt{where} reports the mapping that holds an address? Why does a 1 MiB block live somewhere else than a 1 KiB block? Why does the last line differ from the second word of the first? Why might a latency-sensitive program care?
\omcode{code/interviews-3/05-operating-systems-concurrency-and-networks/cpp/malloc_mmap.cpp}{20}{30}{Where malloc gets its memory. A helper, \texttt{where}, is not shown.\protect\hfil\protect\break}\label{lst:iv3:os:mallocmmap}
\end{question}

\begin{question}[$\star\star$ \iqroles{QD, Dev, LLE} \iqfirm{any} \iqlevel{junior, experienced}]\label{iq:iv3:operating-systems-concurrency-and-networks:12}
\Cref{lst:iv3:os:rssvsz} reports the virtual size and the resident size, in MiB, from \texttt{/proc/self/status}. About what two lines does it print? How can \texttt{mmap} of 1 GiB succeed without using a gigabyte? What happens if every process touches all it asked for?
\omcode{code/interviews-3/05-operating-systems-concurrency-and-networks/python/rss_vsz.py}{1}{23}{Asking for memory and using it.\protect\hfil\protect\break}\label{lst:iv3:os:rssvsz}
\end{question}

\begin{question}[$\star\star$ \iqroles{QD, Dev, LLE} \iqfirm{any} \iqlevel{junior}]\label{iq:iv3:operating-systems-concurrency-and-networks:13}
\Cref{lst:iv3:os:stackdepth} is built with g++ 11.4, \texttt{-O2}, and run with \texttt{ulimit -s 8192}, so the stack limit is 8 MiB. Each call uses a little over 1 KiB of stack. What happens with the argument 1000, and with 100\,000? What about 100\,000 under \texttt{ulimit -s 262144}? Can you catch this failure with \texttt{try} and \texttt{catch}? How big is a new thread's stack under the first limit?
\omcode{code/interviews-3/05-operating-systems-concurrency-and-networks/cpp/stack_depth.cpp}{1}{12}{Going down the stack.\protect\hfil\protect\break}\label{lst:iv3:os:stackdepth}
\end{question}

\begin{question}[$\star\star$ \iqroles{QD, Dev, LLE} \iqfirm{crypto, market maker, hedge fund} \iqlevel{junior, experienced}]\label{iq:iv3:operating-systems-concurrency-and-networks:14}
In \cref{lst:iv3:os:pingpong}, two processes pinned to one core pass a byte back and forth 1000 times through pipes. What is the smallest possible value it can print, and why? What is a context switch, and what does it cost beyond the time of the switch itself? Why does the same ping-pong on two cores with spinning avoid this?
\omcode{code/interviews-3/05-operating-systems-concurrency-and-networks/python/ping_pong.py}{1}{27}{Counting switches.\protect\hfil\protect\break}\label{lst:iv3:os:pingpong}
\end{question}

\begin{question}[$\star\star$ \iqroles{QD, Dev} \iqfirm{any} \iqlevel{junior}]\label{iq:iv3:operating-systems-concurrency-and-networks:15}
Three tasks arrive at time 0 in the order A, B, C, with CPU needs of 8, 4 and 1. Find each one's completion time and the average under round robin with a quantum of 2, then under first-come-first-served, then under shortest-job-first. Ignore the cost of switching. Write round robin in Python. What does it become as the quantum grows? Why not always pick the shortest job?
\end{question}

\begin{question}[$\star\star$ \iqroles{QD, Dev} \iqfirm{any} \iqlevel{junior, experienced}]\label{iq:iv3:operating-systems-concurrency-and-networks:16}
You are given which thread holds each lock, and which lock each thread is blocked on (at most one). In Python, return a list of threads forming a deadlock, or \texttt{None}. Which four conditions must all hold for a deadlock, and which is usually removed in practice? Why can this be done without a general cycle search?
\end{question}

\begin{question}[$\star\star$ \iqroles{QD, Dev, LLE} \iqfirm{any} \iqlevel{junior, experienced}]\label{iq:iv3:operating-systems-concurrency-and-networks:17}
On Linux, what two lines does \cref{lst:iv3:os:mmapfile} print? Which of these results does the manual promise, and which is only what happens here? When would you read a file through a mapping rather than with \texttt{read}? When does a shared mapping's write reach the disk?
\omcode{code/interviews-3/05-operating-systems-concurrency-and-networks/python/mmap_file.py}{1}{15}{A shared and a private mapping of one file.\protect\hfil\protect\break}\label{lst:iv3:os:mmapfile}
\end{question}

\begin{question}[$\star\star$ \iqroles{QD, Dev} \iqfirm{any} \iqlevel{junior, experienced}]\label{iq:iv3:operating-systems-concurrency-and-networks:18}
A process saves its state by writing over one file, and \texttt{write} returned success when the power went. What can be on disk? In Python on Linux, write \texttt{atomic\_save(path, data)} so that a crash leaves either the whole old file or the whole new one. Explain each system call. What remains unprotected, even so?
\end{question}

\begin{question}[$\star\star$ \iqroles{QD, Dev, LLE} \iqfirm{crypto, market maker, hedge fund} \iqlevel{junior, experienced}]\label{iq:iv3:operating-systems-concurrency-and-networks:19}
\Cref{lst:iv3:os:clocks} prints the resolution of \texttt{CLOCK\_MONOTONIC}, and how many times it went backwards in 200\,000 reads. What are the two lines on this machine? You must time a request and also stamp an order with the time of day. Which clock do you use for each, and why? How can a clock read take no system call? Why is resolution not the same as accuracy?
\omcode{code/interviews-3/05-operating-systems-concurrency-and-networks/python/clocks.py}{1}{10}{A monotonic clock.\protect\hfil\protect\break}\label{lst:iv3:os:clocks}
\end{question}

\begin{question}[$\star\star$ \iqroles{QD, Dev, LLE} \iqfirm{any} \iqlevel{junior}]\label{iq:iv3:operating-systems-concurrency-and-networks:20}
What two lines does \cref{lst:iv3:os:acceptlater} print, in which order? The client finishes \texttt{connect} and sends data before the server has called \texttt{accept}. How? Describe the handshake, and what the backlog argument controls. What happens to a new client once the queue is full?
\omcode{code/interviews-3/05-operating-systems-concurrency-and-networks/python/accept_later.py}{1}{8}{Connecting before accept.\protect\hfil\protect\break}\label{lst:iv3:os:acceptlater}
\end{question}

\begin{question}[$\star\star$ \iqroles{QD, Dev, LLE} \iqfirm{any} \iqlevel{junior, experienced}]\label{iq:iv3:operating-systems-concurrency-and-networks:21}
Your server restarts and fails with \texttt{Address already in use}, though no process holds the port. What three lines does \cref{lst:iv3:os:timewait} print on Linux? Which side enters \texttt{TIME\_WAIT}, and why does that state exist? Why does the second line not say \texttt{bound}, though the new socket asked for \texttt{SO\_REUSEADDR}?
\omcode{code/interviews-3/05-operating-systems-concurrency-and-networks/python/time_wait.py}{1}{34}{Restarting a server. State 06 is \texttt{TIME\_WAIT}.\protect\hfil\protect\break}\label{lst:iv3:os:timewait}
\end{question}

\begin{question}[$\star\star$ \iqroles{QD, Dev} \iqfirm{any} \iqlevel{junior, experienced}]\label{iq:iv3:operating-systems-concurrency-and-networks:22}
A 10 Gb/s link has a round-trip time of 1 ms, and nothing is lost. What is the bandwidth-delay product? With a window of 65\,535 bytes, what throughput do you get, and how does window scaling repair that? What is the difference between flow control and congestion control? Then state what happens to the rate as the round-trip time doubles, for a fixed window.
\end{question}

\begin{question}[$\star\star$ \iqroles{QD, Dev, LLE} \iqfirm{any} \iqlevel{junior, experienced}]\label{iq:iv3:operating-systems-concurrency-and-networks:23}
Messages are sent over TCP with a 2-byte big-endian length, and then the body. In Python, write a \texttt{Deframer} whose \texttt{feed(data)} takes bytes as they arrive from \texttt{recv} and returns the complete message bodies. Reject a length above 4096. Why does \texttt{msg = sock.recv(4096)} pass a test on loopback and then fail in production? What about the sending side?
\end{question}

\begin{question}[$\star\star$ \iqroles{QD, Dev, LLE} \iqfirm{crypto, market maker, hedge fund} \iqlevel{junior}]\label{iq:iv3:operating-systems-concurrency-and-networks:24}
On Linux, what three lines does \cref{lst:iv3:os:udptrunc} print? What happened to the bytes \texttt{efgh}? Where does the number 65\,507 come from? How large should a real feed's datagrams be, and why?
\omcode{code/interviews-3/05-operating-systems-concurrency-and-networks/python/udp_trunc.py}{1}{14}{Datagrams and buffers.\protect\hfil\protect\break}\label{lst:iv3:os:udptrunc}
\end{question}

\begin{question}[$\star\star$ \iqroles{QD, Dev, LLE} \iqfirm{crypto, market maker, hedge fund} \iqlevel{junior, experienced}]\label{iq:iv3:operating-systems-concurrency-and-networks:25}
On Linux, what four lines does \cref{lst:iv3:os:udpdrop} print? Why does the option return a different value from the one set? The sender saw no error, so who lost the data, and how can you tell? What would you do for a feed handler that drops during bursts?
\omcode{code/interviews-3/05-operating-systems-concurrency-and-networks/python/udp_drop.py}{1}{22}{A receiver that does not read.\protect\hfil\protect\break}\label{lst:iv3:os:udpdrop}
\end{question}

\begin{question}[$\star\star$ \iqroles{QD, Dev, LLE} \iqfirm{crypto, market maker, hedge fund} \iqlevel{junior, experienced}]\label{iq:iv3:operating-systems-concurrency-and-networks:26}
What is an IPv4 multicast group, and how does a host join one? Run \cref{lst:iv3:os:mcast} with \texttt{join}, then with anything else: what does each print on this Linux? Which Ethernet address does group 239.1.2.3 map to, and which other groups share it? What do \texttt{IP\_MULTICAST\_TTL} and \texttt{IP\_MULTICAST\_LOOP} default to, and why does it matter when a test receives data on one machine but not across two?
\omcode{code/interviews-3/05-operating-systems-concurrency-and-networks/python/mcast_join.py}{1}{20}{Joining a group, on the loopback interface.\protect\hfil\protect\break}\label{lst:iv3:os:mcast}
\end{question}

\begin{question}[$\star\star$ \iqroles{QD, Dev, LLE} \iqfirm{any} \iqlevel{junior, experienced}]\label{iq:iv3:operating-systems-concurrency-and-networks:27}
What two lines does \cref{lst:iv3:os:nonblock} print, where each is what \texttt{connect\_ex} returned and then what \texttt{SO\_ERROR} held? Why is the error on a refused connection not in the first position? Why must you check \texttt{SO\_ERROR} when \texttt{select} says the socket is writable?
\omcode{code/interviews-3/05-operating-systems-concurrency-and-networks/python/nonblock.py}{1}{20}{A connection that does not wait.\protect\hfil\protect\break}\label{lst:iv3:os:nonblock}
\end{question}

\begin{question}[$\star\star$ \iqroles{QD, Dev, LLE} \iqfirm{any} \iqlevel{junior, experienced}]\label{iq:iv3:operating-systems-concurrency-and-networks:28}
What happens to \cref{lst:iv3:os:sigpipe} run with no arguments, and with one? In \cref{lst:iv3:os:firstwrite}, the peer has already closed a TCP connection: why does the first send succeed, and what does the second do? How should a server defend itself against this?
\omcode{code/interviews-3/05-operating-systems-concurrency-and-networks/cpp/sigpipe.cpp}{7}{14}{Writing to a closed peer, in C++.\protect\hfil\protect\break}\label{lst:iv3:os:sigpipe}
\omcode{code/interviews-3/05-operating-systems-concurrency-and-networks/python/first_write.py}{1}{15}{Writing to a closed peer over TCP, in Python.\protect\hfil\protect\break}\label{lst:iv3:os:firstwrite}
\end{question}

\begin{question}[$\star\star\star$ \iqroles{QD, Dev, LLE} \iqfirm{any} \iqlevel{experienced}]\label{iq:iv3:operating-systems-concurrency-and-networks:29}
One CPU, strict priorities. A low-priority task takes a lock at time 0 and holds it for 4 ticks of its own CPU time. A high-priority task arrives at time 1 and needs the lock for 1 tick and then runs 1 more. A middle-priority task arrives at time 2 and needs 20 ticks, with no lock. When does the high-priority task finish? What is this called? What does priority inheritance change, and what is the finish time with it? Write a tick-by-tick simulation in Python, and show how the delay depends on the middle task.
\end{question}

\begin{question}[$\star\star\star$ \iqroles{QD, Dev, LLE} \iqfirm{any} \iqlevel{experienced}]\label{iq:iv3:operating-systems-concurrency-and-networks:30}
In C++20, using one \texttt{std::mutex} and condition variables, write a reader-writer lock in which a waiting writer blocks new readers. Include \texttt{try\_lock\_shared}. Why that preference, and what does it risk? How do you test that readers do not overtake a waiting writer, without relying on timing? Why might a reader-writer lock be no faster than a mutex for short critical sections?
\end{question}

\begin{question}[$\star\star\star$ \iqroles{QD, Dev, LLE} \iqfirm{crypto, market maker, hedge fund} \iqlevel{experienced}]\label{iq:iv3:operating-systems-concurrency-and-networks:31}
In C++20 on Linux, write a mutex on one \texttt{std::atomic<int>} and the \texttt{futex} system call, so that a lock nobody else wants makes no system call at all. Why do you need three states, and not two? What goes wrong if the state goes back to 1 after a sleeper wakes up? Why must \texttt{FUTEX\_WAIT} take the expected value as an argument?
\end{question}

\begin{question}[$\star\star\star$ \iqroles{QD, Dev, LLE} \iqfirm{crypto, market maker, hedge fund} \iqlevel{junior, experienced}]\label{iq:iv3:operating-systems-concurrency-and-networks:32}
A client makes two small writes on a TCP connection, a header and a body, and then waits for the reply. The round-trip time is 0.1 ms, and the server holds an acknowledgement for up to 40 ms unless it has data to send, and it replies only once it has both pieces. Nagle's algorithm is on. What is the time until the reply arrives? Explain the mechanism. Give three ways to fix it, and say which is best. Model it in Python.
\end{question}

\begin{question}[$\star\star\star$ \iqroles{LLE, QD, Dev} \iqfirm{crypto, market maker, hedge fund} \iqlevel{experienced}]\label{iq:iv3:operating-systems-concurrency-and-networks:33}
A UDP market data packet arrives at the network card of a Linux server, and your process is blocked in \texttt{recv}. Describe the path, step by step, until \texttt{recv} returns. At which steps can the packet wait, be dropped, or be copied? Which steps does kernel bypass remove? Where would you look to find drops at each stage?
\end{question}

\begin{question}[$\star\star\star$ \iqroles{QD, Dev, LLE} \iqfirm{any} \iqlevel{experienced}]\label{iq:iv3:operating-systems-concurrency-and-networks:34}
What two lines does \cref{lst:iv3:os:epoll} print? Explain each number. Why does edge-triggered mode demand a nonblocking socket and reading until \texttt{EAGAIN}? Why is \texttt{epoll} better than \texttt{select} or \texttt{poll} for ten thousand mostly idle connections, and when is that no real gain?
\omcode{code/interviews-3/05-operating-systems-concurrency-and-networks/python/epoll_modes.py}{1}{21}{Level and edge.\protect\hfil\protect\break}\label{lst:iv3:os:epoll}
\end{question}

\begin{question}[$\star\star\star$ \iqroles{QD, Dev, LLE} \iqfirm{any} \iqlevel{experienced}]\label{iq:iv3:operating-systems-concurrency-and-networks:35}
With g++ 11.4 on Linux, what three lines does \cref{lst:iv3:os:forklock} print? Explain the child's two values. Suppose the other thread held the lock inside a logging library, and the child now calls that library: what happens? What may a child of a multi-threaded process safely do before \texttt{exec}?
\omcode{code/interviews-3/05-operating-systems-concurrency-and-networks/cpp/fork_lock.cpp}{9}{38}{Fork while a lock is held.\protect\hfil\protect\break}\label{lst:iv3:os:forklock}
\end{question}

\begin{question}[$\star\star\star$ \iqroles{QD, Dev, LLE} \iqfirm{crypto, market maker, hedge fund} \iqlevel{experienced}]\label{iq:iv3:operating-systems-concurrency-and-networks:36}
Two processes share a \texttt{pthread\_mutex\_t} in shared memory, and one dies holding it. \Cref{lst:iv3:os:robust} prints the result of the other's lock: what does it print with no arguments, and with one? Why is the second case so dangerous without a time limit? What must the new owner do on \texttt{EOWNERDEAD}, and why can the library not do that for it?
\omcode{code/interviews-3/05-operating-systems-concurrency-and-networks/cpp/robust.cpp}{9}{39}{A lock whose owner dies.\protect\hfil\protect\break}\label{lst:iv3:os:robust}
\end{question}

\begin{question}[$\star\star\star$ \iqroles{QD, Dev, LLE} \iqfirm{crypto, market maker, hedge fund} \iqlevel{junior, experienced}]\label{iq:iv3:operating-systems-concurrency-and-networks:37}
A feed sends the same packets, numbered from 1, on two UDP lines, A and B, so a packet lost on one can arrive on the other. In Python, write an \texttt{Arbiter} whose \texttt{on\_packet(seq)} takes each arrival, from either line, and returns the numbers now deliverable, in order, each exactly once. If more than \texttt{hold} later packets are waiting behind a missing one, declare a gap and go on. Why not wait forever? How do you choose \texttt{hold}?
\end{question}

\begin{question}[$\star\star\star\star$ \iqroles{QD, Dev, LLE} \iqfirm{crypto, market maker, hedge fund} \iqlevel{experienced}]\label{iq:iv3:operating-systems-concurrency-and-networks:38}
In Python, with \texttt{select.epoll} and one thread, write an echo server for messages framed with a 2-byte length. It must survive messages split at any byte, a client that writes faster than it reads, and a client that closes while replies are pending. Where do pending replies go? When do you ask for \texttt{EPOLLOUT}, and why not always? What does one slow client do to the others? Assume you already have \texttt{Deframer}, whose \texttt{feed(data)} returns a list of complete bodies and raises \texttt{ValueError} on a length above 4096, and \texttt{frame(body)}, which returns the length and the body.
\end{question}

\begin{question}[$\star\star\star\star$ \iqroles{QD, Dev} \iqfirm{crypto, market maker, hedge fund} \iqlevel{experienced}]\label{iq:iv3:operating-systems-concurrency-and-networks:39}
Build reliable in-order delivery on an unreliable link in Python, as a round-by-round model: selective repeat, with a window of $w$ packets. The link may lose, duplicate or reorder packets, and may lose acknowledgements. Why must the receiver acknowledge a duplicate? Why is the window needed at all? Prove it delivers in order and exactly once. How do you test it, including at a loss rate of 1? Compare with TCP.
\end{question}

\begin{question}[$\star\star\star\star$ \iqroles{QD, Dev, LLE} \iqfirm{crypto, market maker, hedge fund} \iqlevel{experienced}]\label{iq:iv3:operating-systems-concurrency-and-networks:40}
A producer process and a consumer process share a ring buffer in \texttt{MAP\_SHARED} memory, with \texttt{std::atomic} indices. In C++20, write it. Why must the indices be lock-free, and how do you check that at compile time? If the producer is killed at an arbitrary instruction, what can the consumer see, and what can it never see? How does a mutex in the same memory compare? ThreadSanitizer does not follow two processes, so how do you test the memory orders?
\end{question}
